\apartado{Contenedores, asincronía, observabilidad y costo}
\subsection{Ejecución y escalamiento}
Se propone empaquetar procesos en contenedores con versiones fijadas. Cloud Run administra instancias y puede escalar un servicio a cero; la primera solicitud posterior puede incorporar tiempo de arranque \cite{cloud-run}. El worker del diseño es un \textbf{servicio HTTP privado invocado por tareas}, no un producto de worker pool asumido equivalente.

Como parámetros iniciales del piloto se plantean mínimo cero instancias, máximo dos para API y una para worker, y concurrencia acotada. Son decisiones ajustables, no mediciones ni límites universales. Deben contrastarse con memoria, tiempos de procesamiento y cola; aumentar instancias sin controlar llamadas a IA también puede aumentar consumo.

\subsection{Procesamiento asíncrono}
Cloud Tasks admite tareas HTTP con autenticación mediante cuenta de servicio \cite{cloud-tasks}. El diseño separa aceptación de solicitud y ejecución. Cada trabajo necesitará estado persistente, reintentos idempotentes, cancelación y recuperación ante interrupción; reintentar una entrega no debe duplicar resultados. Esta integración pertenece a hitos posteriores, aunque H0 define sus contratos.

\subsection{Señales operativas y privacidad}
Se prevén registros estructurados, compatibles con el tratamiento de campos de Cloud Logging \cite{logging}. Un \codigo{requestId} y un \codigo{jobId} permitirían relacionar fallos sin guardar apuntes ni chats. Las señales incluirán latencia, errores por dependencia, antigüedad de cola, fallos de extracción y consumo estimado. Nunca se requieren tokens, URLs firmadas o prompts completos para explicar que una petición falló.

\subsection{Plan económico sin créditos trial}
El Free Tier tiene condiciones y límites; requiere cuenta de facturación, y los excedentes de una cuenta de pago pueden generar cargos \cite{free-tier}. El piloto no usa el crédito de prueba como base de viabilidad. Se propone revisar región, cuotas y consumo de cada servicio antes de crear recursos. \codigo{us-central1} es una región candidata del plan, aún sujeta a residencia de datos y latencia.

\small
\begin{tabularx}{\linewidth}{P{3cm}Y}
\toprule
\textbf{Grupo de gasto} & \textbf{Control y evidencia previstos}\\
\midrule
Run, builds y registro & Tiempo/memoria, instancias, número de imágenes y limpieza.\\
Firestore y Storage & Operaciones, índices, bytes, transferencia y retención física.\\
Tasks y operación & Despachos, reintentos, logs, Scheduler y secretos.\\
IA externa & Tokens/imagen, cuota por cuenta, reserva de consumo y apagado operativo.\\
\bottomrule
\end{tabularx}
\normalsize
\textbf{Presupuesto no equivale a control automático por defecto.} Las alertas por sí mismas no detienen gasto. La documentación vigente distingue presupuestos solo de alerta y presupuestos con límite de gasto; su disponibilidad y cobertura deben comprobarse antes de confiar en ellos \cite{budgets}. No hay costo observado ni monto mensual demostrado en este corte.
